% Generated draft: requires mathematical and source review.
\documentclass[11pt]{article}
\usepackage[margin=25mm]{geometry}
\usepackage{fontspec}
\setmainfont{Latin Modern Roman}
\usepackage{amsmath,amssymb,mathtools,unicode-math}
\setmathfont{Latin Modern Math}
\usepackage{enumitem,needspace,xurl,hyperref}
\setlength{\parindent}{0pt}
\setlength{\parskip}{5pt}
\setlength{\emergencystretch}{3em}
\newcommand{\N}{\mathbb{N}}
\newcommand{\Z}{\mathbb{Z}}
\newcommand{\Q}{\mathbb{Q}}
\newcommand{\R}{\mathbb{R}}
\newcommand{\C}{\mathbb{C}}
\newcommand{\F}{\mathbb{F}}
\newcommand{\A}{\mathbb{A}}
\newcommand{\PP}{\mathbb{P}}
\newcommand{\E}{\mathbb{E}}
\newcommand{\eps}{\varepsilon}
\newcommand{\dd}{\,\mathrm{d}}
\begin{document}
\section*{UnsolvedMath ID: 11100014}
\textbf{Source problem code:} catalogue entry
\subsection{Definitions and mathematical statement}
Fix the prime and height in the source's Morava-theory setting. Let \(K(n)\) denote Morava \(K\)-theory, \(L_{K(n)}\) Bousfield localization at that theory, \(E\) the source's Morava-type \(E\)-spectrum and \(S\) its profinite acting group. The record does not specify whether \(S\) is the stabilizer or an extended stabilizer, nor the coefficient/action conventions; these choices affect descent.

Hovey's Morava K- and E-theory item 4 asks for a more natural construction of the spectral sequence relating continuous cohomology to localized homotopy. In the usual descent grading its intended form is
\[
 H_c^s\!\left(S;\pi_t L_{K(n)}(E\wedge X)\right)
       \Longrightarrow \pi_{t-s}L_{K(n)}X.
\]
The topology on the coefficient module, continuity, class of spectra \(X\), convergence and precise group action are not supplied, so this display records the stated source/target relationship rather than asserting a general convergence theorem.

The source criticizes an Adams--Novikov construction described as starting with \(E\wedge X\wedge F\), for a type-\(n\) object, and then taking an inverse limit. Conventionally a finite \(p\)-local spectrum is type \(n\) when its lower Morava \(K\)-homologies vanish but its \(K(n)\)-homology is nonzero. The actual tower and comparison maps are absent. Hovey says the stated method only works for dualizable \(X\), meaning dualizable in the \(K(n)\)-local category, and seeks a construction less dependent on arbitrary choices. This is a historical limitation report, not a present-day impossibility assertion.

He raises the same naturality objection to the more general Devinatz--Hopkins fixed-point-spectrum construction designed to give continuous rather than discrete cohomology. The request preserves that distinction and asks for an intrinsic conceptual construction; it gives no formal uniqueness, functoriality or canonicity criterion. The supplied partial-progress summary says rigorous constructions and modern frameworks exist, but does not provide an exact theorem, references or a completed answer to the aesthetic/naturality requirement.
\subsection{Short English statement}
Give a more intrinsic construction of the continuous-cohomology Morava descent spectral sequence and the related Devinatz--Hopkins fixed-point spectra. The source objects to inverse-limit choices and a dualizability restriction. Its desired notion of naturality and the precise descent hypotheses remain unspecified.
\subsection{Sources}
\begin{itemize}
\item Local UnsolvedMath registry, permanent ID \(11100014\), exact source code \texttt{catalogue entry}, ``Morava K- and E-theory, item 4,'' attributed to Mark Hovey; no proposed year supplied.
\item \emph{Hovey's Algebraic Topology Problem List}, surviving mirror: \url{https://www-users.cse.umn.edu/~tlawson/hovey/}. Background: 65 top-level items across nine category pages; HTML-ordered-list extraction, accessed 2026-07-29.
\item The source names Adams--Novikov spectral sequences and Devinatz--Hopkins fixed-point constructions without complete references. The research summary refers to rigorous modern formulations and greater naturality but gives no specific bibliography.
\item Source evidence NEEDS\_REVIEW presents an open problem and calls for a current primary-literature audit. Registry status partially\_solved and classification PARTIAL-PROGRESS describe the remaining intrinsic/canonical-construction issue as not a closed specific theorem.
\item Dataset attribution: Ulam AI / UnsolvedMath Contributors, CC BY 4.0. Source rights NEEDS\_REVIEW: publicly accessible author/publisher copy, redistribution terms requiring review.
\end{itemize}
\end{document}
